\fichatitulo{03 · CONTEXTO INSTITUCIONAL}{Relatório · 2024}{Use cases for AI in parliaments}
{\small INTER-PARLIAMENTARY UNION.
\textit{Use cases for AI in parliaments}. [S. l.]: IPU, dez. 2024.
\href{https://www.ipu.org/file/20635/download}{Documento institucional disponível no portal da IPU}. Acesso em: 6 set. 2026.\par}

\campo{Problema e objetivo}
Apoiar parlamentos na definição de aplicações de IA, organizando necessidades institucionais em descrições de funcionamento e critérios para acompanhar sua implementação.

\campo{Principal contribuição · síntese interpretativa}
Reúne casos de uso elaborados por parlamentos em uma estrutura compartilhável. Distingue a especificação de como um sistema deve funcionar do relato de uma implementação efetivamente realizada (Introdução, p. 3).

\campo{Estrutura e recorte consultado}
Leitura da introdução e do caso 024, sobre consultas em linguagem natural ao portal do Senado italiano. Esse caso descreve atores, pré-requisitos, fluxo, alternativas, dados, integrações e métricas. Inclui tempo de resposta, satisfação e relevância dos resultados (p. 150--151).

\campo{Resumo extrativo · fragmentos literais selecionados}
\begin{itemize}
\excerto{Natureza do documento}{A use case describes how a system should work.}{p. 3, Introdução; página 4 do PDF.}
\excerto{Qualidade da busca}{Accuracy and relevance of search results}{p. 151, caso 024, Success metrics; página 152 do PDF.}
\excerto{Experiência do usuário}{User satisfaction ratings and feedback}{p. 151, caso 024, Success metrics; página 152 do PDF.}
\end{itemize}

\campo{Limitações e alcance}
O recorte descreve objetivos e resultados esperados, sem apresentar mensurações de desempenho. As métricas listadas não vêm acompanhadas de um protocolo completo de avaliação. Esta ficha não sintetiza todos os casos da coletânea.

\campo{Relação com a dissertação · análise proposta}
Ajuda a explicar por que avaliar um RAG legislativo interessa à administração pública. Pode orientar a discussão sobre finalidade e utilidade, mas não comprova eficácia. Para nossa pesquisa, a questão é quais expectativas institucionais se tornam observações mensuráveis e quais ficam fora do protocolo. Satisfação e tempo de resposta não devem ser prometidos como resultados se não houver instrumentos e execução correspondentes.

\registro{Escopo consultado nesta preparação: Introdução (p. 3) e caso 024 (p. 150--151), edição de dezembro de 2024. Leitura parcial da coletânea. Chave no projeto: \texttt{ipu2024aiParliaments}. Excertos em inglês, sem tradução.}
